\documentclass[../main.tex]{subfiles}
\begin{document}
%==========================================TIÊU ĐỀ TẬP========================================
\begin{table}[h]
    \begin{adjustwidth}{-1cm}{}
        \begin{tabular}{>{\centering\arraybackslash}p{6cm}|>{\raggedright\arraybackslash}p{12.5cm}}
            \multirow{5}{6cm}
            {\\[-40pt]
                \flaregothic\fontsize{20pt}{20pt}\selectfont \centering
                \color{eptype}HISTOIRE\color{black}\\
                \freshbot\fontsize{72pt}{72pt}\selectfont
                \color{epnum}43\\[6pt]
                \cafeteria\fontsize{20pt}{20pt}\selectfont\color{epnum}(D-43)\color{black}
                \\[18pt]
                \color{black}
            }
             &
            % Dòng thứ nhất: tên tập tiếng Nhật
            {\fotskip\fontsize{18pt}{18pt}\selectfont \ruby{新}{あたら}しい\ruby{日々}{ひび}}
            \\[6pt]
            % Dòng thứ hai: tên tập tiếng Anh
             & {\cafeteria\fontsize{18pt}{18pt}\selectfont New Days}                                                                                                                                           \\[4pt]
            % Dòng thứ ba: tên tập tiếng Việt
             & {\vnmsans\fontsize{18pt}{18pt}\selectfont Những ngày đầu tiên}                                                                                                                                  \\[8pt]
            % Dòng thứ tư: Nhân vật xuất hiện
             & {\caviar{\fontsize{13pt}{13pt}\selectfont Nhân vật: \emp{kuon2} \emp{kaiho2} \emp{ryuuhou2} \emp{suzuno2} \emp{naomi2} \emp{game} \emp{delu} \emp{kson} \emp{mike1} \emp{koneko} \emp{miyako}}}
            \\[8pt]
        \end{tabular}
    \end{adjustwidth}
\end{table}
%=========================================TRUYỆN CHÍNH========================================
\color{black}

\txtbx[2]{
    {\color{vol} \uline{Tóm tắt:}} Sau đêm tái ngộ kỳ diệu mở ra cuộc sống mới dưới nhân dạng con người, ba mẹ con Miyako, Toa và Kiri chính thức bắt đầu hành trình hòa nhập vào đại gia đình Hikawa. Với ba bộ đồng phục hầu gái đặc chế vừa vặn cho đôi tai và chiếc đuôi cún, mỗi người được Cậu chủ Kuon giao phó một trọng trách riêng biệt. Những tình huống dở khóc dở cười liên tục nổ ra: từ cuộc họp bảo mật triệt để với các cư dân trong ký túc xá, niềm vui khám phá ẩm thực và những sở thích, sở đoản mới mẻ khi hóa người, sức khỏe kinh ngạc của Toa ở xưởng cơ khí và phòng lab, những câu hỏi truy vấn logic của Kiri bên bàn học Naomi, cho đến màn chạm trán tóe lửa giữa tộc cẩu và cặp chị em miêu nữ Mikeneko - Koneko. Đằng sau sự nhộn nhịp ấy là một mái ấm chan chứa tình thân, nơi ngọn lửa sinh mệnh của những người bạn bốn chân năm nào được thắp sáng lại rực rỡ hơn bao giờ hết.
}

\pov{Hikawa Kuon}

Việc đón nhận ba thực thể sinh học mang hình hài con người nhưng sở hữu linh hồn và tập tính của loài chó bước vào nhịp sống thường nhật của ký túc xá phố Anwa hóa ra lại là một thử thách kỳ lạ vượt xa mọi dự liệu của tôi. Ừ thì, đúng là ba người ấy đã từng là ba con chó trong nhà Hikawa chúng tôi, nhưng chúng tôi vẫn không tránh khỏi cảm giác gì đó lạ lẫm.

Kể từ buổi rạng sáng hôm đó, không khí ảm đạm bao trùm căn nhà suốt hơn một tháng ròng sau vụ trộm mộ bỗng chốc tan biến hoàn toàn. Thay vào đó là một chuỗi những xáo trộn vừa ngỡ ngàng, vừa khiến người ta không khỏi bật cười.

Dù đã mang vóc dáng của những thiếu nữ hoàn chỉnh với trí tuệ và khả năng ngôn ngữ lưu loát thừa hưởng từ công nghệ biến đổi gen của Koyori, những thói quen và bản năng loài chó vốn ăn sâu vào tiềm thức suốt bao năm qua vẫn chẳng thể nào lập tức biến mất trong một sớm một chiều.

Khứu giác và thính giác của cả ba người nhạy bén đến mức khó tin. Mỗi buổi chiều, khi chiếc xe máy phân khối lớn của bố tôi chỉ vừa rẽ vào đầu con phố cách nhà tới gần năm trăm mét, đôi tai cún màu trắng muốt trên đầu Miyako đã khẽ giật giật hai nhịp. Cô lập tức đặt chiếc giẻ lau xuống bàn, nhẹ nhàng bước xuống sảnh tầng một, cẩn thận xếp sẵn đôi dép lê đi trong nhà ngay ngắn trước cửa cuốn để chờ đón ông ấy trở về.

Toa thì sở hữu nguồn năng lượng dồi dào đến mức gần như không bao giờ chịu ngồi yên một chỗ. Mỗi lần nghe thấy tiếng huýt sáo vô tình của ai đó hay tiếng chuông báo tin nhắn điện thoại, chiếc đuôi xù dài màu trắng sau lưng cô bé lại ngoáy tít mù thành những vòng tròn hớn hở, hai tai vểnh cao hết cỡ, đôi mắt màu vàng sáng bừng lên vẻ phấn khích tột độ. Trong khi đó, Kiri lại có thói quen tìm đến những khoảng không gian hẹp và tĩnh lặng, lặng lẽ ngồi xếp bằng nơi góc phòng khách để quan sát từng cử chỉ, từng thói quen của mọi người trong nhà với ánh mắt vàng kim sắc sảo và trầm ngâm.

Và buồn cười nhất là những lúc tôi vô thức làm theo thói quen cũ. Đôi khi đi ngang qua hành lang, thấy Toa hay Kiri đang đứng ngẩn ngơ chờ chỉ thị, bàn tay tôi lại tiện thể đưa lên xoa nhẹ vào khoảng giữa hai tai của chúng y hệt như những ngày còn ở xưởng cơ khí ngoại ô.

Những lúc ấy, cả hai em đều theo phản xạ nhắm tịt mắt lại, phát ra những tiếng hừ hừ trong cổ họng đầy thỏa mãn, chiếc đuôi đung đưa lia lịa. Thế rồi, chỉ một giây sau đó, nhận ra bản thân giờ đây đã là những thiếu nữ mang thân xác con người, gương mặt của Toa và Kiri bỗng đỏ bừng lên như quả cà chua chín. Cả hai luống cuống lùi lại nửa bước, hai tay ôm lấy má, đôi tai cún cụp hẳn xuống vì ngượng ngùng khiến tôi cũng phải gãi đầu bối rối.

Việc hóa thành người cũng mở ra cho ba mẹ con một chân trời hoàn toàn mới mẻ, mà điều kỳ diệu nhất đối với họ chính là thế giới ẩm thực.

Thuở trước, khi còn là những chú chó bốn chân, có vô vàn món ăn bị liệt vào danh sách cấm kỵ tuyệt đối: sô-cô-la, nho, hành tây, tỏi, các món sốt đậm vị hay những bát mì ramen bốc khói ngập tràn hành hoa - thực tế là chúng có thể ăn được, chỉ là không được phép ăn thường xuyên, và sô-cô-la chúng tôi ăn là loại sô-cô-la sữa. Giờ đây, với hệ tiêu hóa và vị giác hoàn thiện của con người, ba mẹ con đã có thể đường hoàng nếm trải trọn vẹn những phong vị mà trước đây họ chỉ có thể hít hà từ xa qua làn khói bếp.

Tôi vẫn nhớ như in buổi trưa hôm thứ hai, khi mẹ tôi mang ra một đĩa bánh ngọt phủ sô-cô-la đen cùng một chùm nho mẫu đơn mát lạnh. Toa cầm thanh bánh lên cắn một miếng nhỏ, đôi mắt vàng của em ấy bỗng mở to hết cỡ, chiếc đuôi trắng ngoáy lia lịa đến mức suýt quạt ngã chiếc ghế bên cạnh vì sung sướng. Kiri thì nhấm nháp từng quả nho một cách từ tốn, hai má ửng hồng, khẽ thì thầm rằng vị ngọt thanh mát này kỳ diệu hơn bất kỳ loại thức ăn đóng hộp nào mà Koyori từng đưa cho họ ở phòng thí nghiệm ngầm.

Song song với việc thích nghi thể xác, mỗi người trong số họ cũng dần bộc lộ những sở thích và cả những sở đoản rất đỗi con người.

Miyako có niềm đam mê mãnh liệt với công việc bếp núc và chăm sóc gia đình. Cô học hỏi các công thức nấu ăn mới từ mẹ tôi với tốc độ đáng kinh ngạc, thích tự tay hái những búp trà thảo mộc phơi khô rồi pha chế thành những tách trà thanh mát giúp mọi người an thần. Thế nhưng, người mẹ cún đoan trang này lại mắc một sở đoản dở khóc dở cười: cô hoàn toàn mù tịt về các thiết bị điện tử hiện đại. Chiếc máy giặt tự động với hàng chục phím cảm ứng, chiếc lò vi sóng điện tử hay chiếc điều khiển điều hòa từ xa đối với Miyako chẳng khác nào những cỗ máy ma thuật phức tạp. Mỗi lần bảng điều khiển phát ra tiếng bíp báo lỗi, Miyako lại giật thót mình, chắp hai tay trước ngực cúi gập người xin lỗi chiếc máy giặt vì sợ bản thân vừa làm hỏng tài sản của gia đình.

Toa thì phóng khoáng, mê vận động, thích ăn thịt nướng và những món đồ ngọt béo ngậy. Cô bé thích thú với việc chạy nhảy, leo trèo và sẵn sàng dùng đôi tay khỏe khoắn để bê vác mọi thứ. Tuy nhiên, sở đoản lớn nhất của Toa chính là sự khéo léo tỉ mỉ. Đôi bàn tay vốn quen nhấc bổng cả tạ sắt thép của cô bé cứ chạm vào kim chỉ, việc gấp góc chăn hay cầm những chiếc đĩa sứ mỏng manh là lại lóng ngóng vụng về, nhiều phen suýt bóp nát chiếc thìa trà vì dùng lực quá tay.

Ngược lại hoàn toàn với chị mình, Kiri yêu thích sự yên tĩnh, trật tự và ngăn nắp. Cô bé say mê đọc những cuốn bách khoa toàn thư, thích giải những câu đố toán học hóc búa, và món khoái khẩu của Kiri là những que kem matcha mát lạnh mang vị đắng nhẹ tinh tế. Thế nhưng, vị quân sư nhỏ tuổi với vẻ mặt lúc nào cũng điềm đạm, lạnh lùng ấy lại có một điểm yếu chí mạng: cô bé cực kỳ sợ các loài côn trùng nhiều chân. Chỉ cần vô tình nhìn thấy một chú nhện nhỏ giăng tơ nơi góc trần nhà hay một chú bọ cánh cứng bay lạc qua cửa sổ, vẻ thông thái của Kiri lập tức bay biến. Đôi tai cún của cô bé cụp chặt sát vào tóc, cả người run rẩy lùi tít ra sau cánh cửa, luống cuống nắm chặt lấy vạt áo tôi hoặc núp sau lưng Toa để cầu cứu.

Tuy nhiên, vấn đề cấp bách đầu tiên nảy sinh trong hai ngày đầu tiên chính là trang phục. Ba mẹ con không thể mặc mãi những bộ quần áo ngủ hay đồ mặc ở nhà mượn tạm của mẹ tôi, Suzuno và Ryuuhou. Khổ người của Toa khỏe khoắn nên mặc áo phông của Ryuuhou thì chật nách, còn váy của Suzuno thì lại quá dài đối với vóc dáng nhỏ nhắn của Kiri. Chưa kể, những bộ trang phục thông thường hoàn toàn không có khoảng hở để đôi tai và chiếc đuôi cún cử động thoải mái.

Chính vì thế, mẹ tôi và Suzuno đã đích thân lấy số đo, liên hệ với một tiệm may đo thủ công quen biết ở ngoại ô để đặt làm riêng ba bộ đồng phục hầu gái đặc chế.

Vào buổi sáng ngày thứ ba, khi ba chiếc hộp các-tông chứa trang phục mới được chuyển tới, cả tầng năm của ký túc xá như biến thành một buổi trình diễn thời trang nhỏ.

Khi ba mẹ con bước ra khỏi phòng thay đồ, cả tôi, bố mẹ và các em gái đều phải sững sờ trước diện mạo mới mẻ và vô cùng đẹp mắt của họ.

Đứng ở vị trí trung tâm là Miyako. Mái tóc dài màu bạch kim của cô buông xõa bồng bềnh qua bờ vai, đôi tai cún trắng muốt khẽ vểnh nhẹ, được điểm xuyết bằng một vòng hoa nhỏ màu vàng nhạt thanh nhã cài khéo léo bên chiếc bờm hầu gái viền bèo tròn trắng tinh. Cô diện chiếc đầm liền thân cổ điển màu đen tuyền tay dài trang trọng, phần cổ tay áo bo lại bằng lớp viền bèo trắng xếp ly tỉ mỉ. Phía trước cổ áo viền ren là chiếc nơ lụa màu vàng đồng được thắt ngay ngắn, phẳng phiu. Chiếc tạp dề trắng tinh khôi phủ kín từ ngực xuống thân váy dài quá đầu gối kín đáo, phần dây tạp dề thắt nơ lớn sau lưng được khoét một khe hẹp chuẩn xác, để lộ chiếc đuôi xù trắng muốt đung đưa nhịp nhàng. Dưới chân, Miyako mang đôi tất trắng ngắn viền bèo ngang mắt cá kết hợp cùng đôi giày búp bê da đen bóng có quai cài. Trên tay cô cầm một chiếc chổi rơm cán tre truyền thống. Dáng đứng đoan trang, một tay đặt nhẹ trước ngực, ánh mắt vàng kim dịu dàng toát lên phong thái mẫu mực của một nữ quản gia trưởng đĩnh đạc, chu đáo và bao dung.

Đứng bên tay trái là Toa, cô hầu gái mang nguồn năng lượng bùng nổ. Mái tóc trắng dài hơi gợn sóng của Toa được để xõa tự do, đôi tai vểnh cao linh hoạt cùng vòng hoa nhỏ cài lệch một bên. Toa khoác lên mình chiếc đầm đen dáng ngắn ngang đùi trẻ trung, gấu váy viền bèo trắng xếp tầng xòe nhẹ theo từng bước chân tung tăng. Chiếc tạp dề trắng ôm gọn lấy vòng eo thon thả, làm nổi bật chiếc nơ cổ màu vàng đồng rực rỡ. Điểm nhấn ấn tượng nhất trên người Toa chính là đôi chân thon dài mang quần tất đen mỏng ôm sát vô cùng khỏe khoắn, phối cùng đôi giày Mary Jane đen bóng. Vừa mặc xong bộ đồ mới, Toa đã hớn hở giang rộng hai tay, nụ cười tươi rói để lộ chiếc răng khểnh tinh nghịch, chiếc đuôi trắng vẫy tít mù phía sau. Thậm chí để chứng minh sự thoải mái của bộ váy, cô bé còn tiện tay dùng một cánh tay nâng bổng ba chồng thùng các-tông phụ tùng cơ khí nặng trịch đặt cạnh cửa lên cao một cách nhẹ tênh như không khí.

Ở phía bên phải, Kiri lại mang một phong thái hoàn toàn đối lập. Mái tóc bạch kim của cô bé được cắt ngắn ngang vai, uốn lọn nhẹ nhàng ôm lấy gương mặt thanh tú và búi gọn gàng sau gáy, bên tai cún cài một đóa hoa nhỏ trang nhã. Ánh mắt màu vàng của Kiri điềm đạm, toát lên vẻ sắc sảo và thông tuệ. Cô mặc chiếc đầm đen dài ngang gối tao nhã, chiếc tạp dề trắng viền bèo thắt nơ vừa vặn quanh eo, nơ cổ vàng đồng thắt chặt chẽ. Đôi chân Kiri mang đôi tất trắng cao dưới đầu gối có viền bèo đính nơ đen nhỏ nhắn, và điểm xuyết ở đùi trái là một chiếc vòng nịt đùi bằng da đen mảnh mai đính nơ vô cùng tinh tế. Hai tay Kiri nâng một chiếc khay kim loại tròn sáng loáng ngang ngực, tư thế đứng thẳng thớm, ánh mắt quan sát mọi người xung quanh với sự bình tĩnh và chuẩn mực của một hầu gái danh gia vọng tộc.

Nhìn ba người đứng ngay ngắn thành một hàng trước mặt, bố tôi không giấu nổi nụ cười hiền hậu, còn mẹ tôi thì tấm tắc khen ngợi sự khéo tay của người thợ may.

Tôi hắng giọng một tiếng, đặt chiếc máy tính bảng xuống bàn trà, chính thức bước vào buổi họp phân công nhiệm vụ: ``Được rồi. Đồ đạc đã tươm tất, từ hôm nay ba người sẽ chính thức tham gia vào các công việc chung của nhà Hikawa. Vì mỗi người có một thế mạnh và tính cách riêng, anh đã bàn bạc với Game để phân chia trách nhiệm cụ thể như sau.''

Miyako, Toa và Kiri lập tức đứng nghiêm trang, chăm chú lắng nghe từng lời của tôi.

``Thứ nhất, Miyako sẽ giữ vai trò Hầu gái trưởng,'' tôi nhìn người phụ nữ tóc trắng. ``Em sẽ phụ trách quản lý tổng thể công việc nội trợ, dọn dẹp các khu vực sinh hoạt chung của ký túc xá, chăm lo bữa ăn hàng ngày cho gia đình cùng mẹ anh. Đồng thời, em sẽ hỗ trợ trực tiếp cho anh với Kaiho trong việc sắp xếp hồ sơ, quản lý các đầu việc hậu cần và lịch trình làm việc cá nhân.''

Miyako khẽ nghiêng người, đôi bàn tay thon thả đặt ngay ngắn trước tạp dề, cúi đầu cung kính: ``Em xin tuân lệnh Cậu chủ. Em nhất định sẽ quán xuyến mọi việc chu đáo để Cậu chủ không phải bận lòng.''

Tôi quay sang Toa, cô bé đang đứng nhấp nhổm với đôi mắt sáng rực chờ đợi: ``Toa, em có thể lực và sức mạnh phi thường. Trọng trách của em là hỗ trợ các công việc nặng nhọc. Buổi sáng, em sẽ phụ bố anh ở xưởng cơ khí, bốc dỡ các khối linh kiện và phụ tùng sắt thép. Những lúc Suzuno làm việc ở phòng thí nghiệm mới xây trên bãi đất hoang, em sẽ sang đó phụ bê vác các bình chứa và thiết bị cồng kềnh. Ngoài ra, những ngày Ryuuhou đi tập nhạc cùng ban nhạc underground, em sẽ làm trợ lý hậu cần, phụ trách mang vác dàn loa đài và âm ly cho con bé. Em thấy sao, có kham nổi không?''

Nghe đến đó, Toa nhảy cẫng lên một nhịp, hai tay nắm chặt lại đầy hào hứng: ``Tuyệt quá đi mất! Cứ việc nặng là giao cho em, em bảo đảm một mình em vác hết cả xưởng xe cũng được luôn ạ!''

Ryuuhou đứng tựa lưng vào tường lập tức huýt sáo một tiếng tán thưởng: ``Được đấy, Toa. Chiều nay có buổi tập dợt, chị sẽ cho em thử sức với dàn loa sub năm mươi cân xem sao nhé.''

Cuối cùng, tôi nhìn sang Kiri: ``Còn Kiri, em có tư duy logic sắc bén và tính tình cẩn thận. Em sẽ phụ trách chăm sóc, bảo vệ và nhắc nhở lịch học tập, sinh hoạt cho Naomi. Bên cạnh đó, em sẽ hỗ trợ kiểm kê dụng cụ kỹ thuật cùng Game, và phụ trách bàn trà tiếp khách ở các tầng khi có người ghé thăm.''

Kiri nâng chiếc khay bạc lên cao hơn một chút, khẽ gật đầu, giọng nói phẳng lặng nhưng rành rọt: ``Em đã ghi nhận đầy đủ. Tỷ lệ sai sót trong việc quản lý thời gian của Naomi-sama sẽ được duy trì ở mức dưới không phẩy một phần trăm.''

Naomi đang ngồi trên ghế sofa nghe thấy tên mình liền ngẩng đầu lên, đôi mắt tròn xoe nhìn Kiri chớp chớp, miệng khẽ nở một nụ cười bẽn lẽn.

Nhìn nét mặt hân hoan của ba mẹ con, tôi dừng lại một chút, khẽ nhíu mày lo lắng rồi hỏi thêm bằng giọng ân cần: ``Nhưng trước khi bắt tay vào làm, anh muốn hỏi thật ba người. Khối lượng công việc anh vừa phân chia như vậy có bị quá nhiều không? Có khiến mọi người cảm thấy bất tiện hay quá sức ở điểm nào không? Nếu có bất kỳ điều gì cảm thấy không thoải mái, các em cứ thẳng thắn nói với anh ngay, chúng ta hoàn toàn có thể sắp xếp và điều chỉnh lại bất cứ lúc nào.''

Nghe câu hỏi đầy quan tâm ấy, Miyako xúc động lắc đầu thật nhẹ: ``Cậu chủ không cần phải bận tâm đâu ạ. So với những bài huấn luyện khắc nghiệt trong bóng tối trước kia, những công việc này ngập tràn ánh sáng và hơi ấm tình thân. Được làm việc cho gia đình là niềm hạnh phúc lớn lao của chúng em, hoàn toàn không có chút bất tiện nào.''

Toa cười toe toét, lắc đầu nguầy nguậy: ``Đúng đó anh Kuon ơi! Em còn thấy việc hơi ít ấy chứ, em chỉ sợ không có đủ đồ nặng cho em vác thôi!''

Kiri cũng đẩy nhẹ chiếc khay bạc, đưa ra nhận định lạnh lùng nhưng chuẩn xác: ``Theo các chỉ số sinh trắc học và mức tiêu thụ calo trung bình, định mức lao động này chỉ chiếm xấp xỉ ba mươi hai phần trăm công suất tối đa của cơ thể biến đổi gen. Đây là một lịch trình vận động hoàn toàn tối ưu và có lợi cho sự phát triển cơ bắp của chúng em ạ.''

Tôi khẽ mỉm cười nhẹ nhõm trước câu trả lời của ba người.

\ct

Tuy nhiên, trước khi chính thức để ba mẹ con xuất hiện công khai tại các tầng sinh hoạt chung, một nhiệm vụ mang tính sống còn khác cần phải được thực hiện triệt để: thiết lập hàng rào bảo mật.

Ngay trong buổi chiều hôm đó, tôi đã triệu tập một cuộc họp riêng đặc biệt tại phòng họp nhỏ tầng năm. Thành phần tham dự chỉ gồm các cư dân không thuộc gia đình Hikawa: Kson từ tầng sáu, Delutaya từ tầng hai, cùng cặp chị em Mikeneko và Koneko từ tầng tám.

Ngồi trước chiếc bàn gỗ dài, tôi nhìn thẳng vào mắt từng người, vẻ mặt nghiêm nghị chưa từng có: ``Hôm nay anh gọi mọi người tới đây là để quán triệt một vấn đề tuyệt mật liên quan trực tiếp đến sự an toàn của toàn bộ ký túc xá này.''

Kson khẽ nhướng mày, đặt lon nước ngọt xuống bàn: ``Gắt thế Kuon-paisen? Liên quan đến ba cô em hầu gái có tai cún mới tới sáng nay à?''

``Chính xác,'' tôi gật đầu, giọng đanh lại. ``Ba người họ là trường hợp đặc biệt liên quan đến thí nghiệm sinh học ngầm của một tổ chức bí mật bên ngoài. Anh yêu cầu tất cả mọi người ở đây: tuyệt đối không được tiết lộ bất kỳ thông tin, hình ảnh hay lời nói nào về sự hiện diện và lai lịch của Miyako, Toa và Kiri ra thế giới bên ngoài. Không được nói với nhân viên ở COVER Corp, không được nhắc tới với các thành viên \hll\ khác, và tuyệt đối cấm chia sẻ trên mạng xã hội hay các buổi phát sóng cá nhân.''

Tôi ngừng lại một nhịp, ánh mắt quét qua từng khuôn mặt: ``Nếu có bất kỳ ai vô tình nhìn thấy họ, mấy đứa phải nói là: họ là những người họ hàng xa ở vùng quê lên phụ giúp việc nhà cho gia đình Hikawa. Tổ chức tạo ra họ có thể đang lùng sục khắp nơi. Chỉ cần một vết rò rỉ nhỏ nhất, hiểm họa sẽ không chỉ ập xuống ba mẹ con họ, mà cả căn nhà này cũng sẽ bị lôi vào vòng nguy hiểm. Anh thực sự nghiêm túc đấy, nên là anh mong mọi người hiểu tính chất sống còn của nó.''

Không gian căn phòng chìm vào sự im lặng trang nghiêm.

Kson khoanh hai tay trước ngực, nét cười cợt thường ngày biến mất hoàn toàn, thay vào đó là khí chất đanh thép của một cựu Tổng trưởng giang hồ: ``Hiểu rồi. Dân trong giới của tôi hiểu rất rõ luật im lặng. Anh đã cưu mang em ở đây thì lời anh nói là mệnh lệnh. Em sẽ bảo đảm tầng sáu sẽ kín như bưng, không một con ruồi nào lọt được thông tin ra ngoài.''

Delutaya khẽ siết chặt vạt áo, đôi mắt to tròn ánh lên vẻ kiên định: ``Em từng nếm trải cảm giác bị săn đuổi và lộ thông tin cá nhân đau đớn thế nào rồi\dots\ Em xin thề sẽ giữ kín miệng, tuyệt đối không bao giờ để ba mẹ con họ rơi vào hoàn cảnh như em trước kia đâu ạ.''

Ở góc bàn, Mikeneko dù đôi lúc hay cằn nhằn nhưng là người từng trải qua những cơn bão truyền thông kinh hoàng nhất. Nàng Mèo hồng ôm chặt lấy em gái Koneko vào lòng, khẽ thở dài một tiếng nhưng ánh mắt lại vô cùng nghiêm túc: ``Anh không cần phải nhắc lại đâu, Kuon-senpai. Em hiểu rõ cái giá của việc bị thế giới ngoài kia nhòm ngó. Em với Koneko sẽ coi như họ chỉ là người làm bình thường của nhà anh, không tò mò, không hé răng nửa lời.''

Koneko ngước nhìn chị gái rồi nhìn tôi, ngoan ngoãn gật đầu đồng tình.

Thấy sự đồng lòng tuyệt đối của mọi người, tảng đá đè nặng trong lòng tôi bấy lâu mới được trút bỏ phần nào. Tôi cúi đầu cảm ơn họ: ``Cảm ơn mọi người. Anh nợ mọi người một ân tình lớn.''

\ct

Những ngày đầu tiên sau buổi phân công nhiệm vụ diễn ra với muôn vàn sắc thái sinh động.

Ở xưởng cơ khí của bố tôi dưới tầng một và khu vực ngoại ô, Toa nhanh chóng trở thành một hiện tượng kỳ lạ khiến cánh thợ cơ khí quen biết phải trố tròn mắt kinh ngạc. Những khối động cơ mô tô bốn xi-lanh nặng gần một tạ vốn phải cần tới hai thanh niên lực lưỡng dùng đòn bẩy khiêng vác thì Toa chỉ cần một tay đã có thể nhấc bổng lên đặt ngay ngắn vào bệ thử nghiệm. Mỗi lần làm xong một việc nặng, cô bé lại quay sang nhìn bố tôi với ánh mắt háo hức mong chờ được khen ngợi, chiếc đuôi cún sau váy đen ngoáy lia lịa khiến bố tôi chỉ biết cười xòa, vỗ vai cô bé khen ngợi rồi dúi vào tay Toa một que kem mát lạnh.

Thế nhưng, sự xông xáo quá mức của Toa cũng tạo ra không ít pha thót tim. Trong một buổi chiều sang phụ Suzuno dọn dẹp phòng thí nghiệm mới, vì quá nhiệt tình bê một thùng bình phản ứng hóa học bằng thủy tinh chịu lực, Toa đã vội vã lao qua cửa mà quên mất phải cúi đầu, khiến mép thùng va mạnh vào khung cửa thép phát ra một tiếng *Keng* chát chúa.

Cô chị Năm đang cặm cụi pha chế bên bàn vi lượng giật bắn mình, hai tay ôm ngực thảng thốt: ``Á\dots\ Toa ơi, bình tĩnh nào! Thùng đó là dung môi este tinh khiết đấy, làm rơi là nổ tung cả gian phòng lab bây giờ!''

Toa đứng chết trân tại chỗ, hai tai cún cụp hẳn xuống đầu, mặt mày tái mét như sắp khóc vì sợ bị đuổi về. May mắn thay, Kiri từ phía sau đã kịp thời bước tới, chiếc khay bạc trên tay đỡ lấy phần đáy thùng một cách chuẩn xác: ``Trọng tâm lệch mười lăm độ về phía tây nam. Hệ số ma sát của găng tay cao su không đủ. Lần sau Toa giảm tốc độ di chuyển xuống bốn mươi phần trăm ạ.''

Suzuno thở phào nhẹ nhõm, nhìn gương mặt ỉu xìu đầy hối lỗi của cô Chó chị rồi lại nhìn nét mặt nghiêm túc của Kiri, cô thiếu nữ mười sáu tuổi khẽ bật cười. Suzuno bước tới, đưa tay xoa nhẹ đầu Toa rồi khẽ vỗ vai Kiri: ``Không sao đâu, chị không mắng em đâu mà. Có hai đứa ở đây giúp đỡ, gian phòng lab này vui vẻ hơn hẳn trước kia đấy.''

Được xoa đầu, đôi tai của Toa lập tức dựng đứng trở lại, nụ cười rạng rỡ lại bừng sáng trên môi.

\ct

Tại phòng khách tầng năm, Kiri cũng nhanh chóng thể hiện tài năng đặc biệt của mình. Trí óc được tôi luyện qua các bài kiểm tra của Koyori kết hợp với bản tính tò mò khiến Kiri trở thành một cỗ máy truy vấn không ngừng nghỉ.

Mỗi buổi tối khi ngồi kèm Naomi học bài, Kiri lại đưa ra hàng loạt câu hỏi khiến bé út nhiều phen bối rối: ``Tại sao Naomi-sama lại đặt dấu ngoặc đơn ở vị trí này? Nếu đổi biến số sang vế bên phải thì thời gian tính toán sẽ rút ngắn được hai giây, tại sao sách giáo khoa lại không làm thế? Còn chiếc bút chì màu hồng này, tại sao lại có mùi dâu tây nhân tạo? Mùi hương này có tác dụng kích thích tế bào thần kinh thị giác không?''

Naomi chớp chớp mắt, ngây ngô đáp: ``Vì\dots\ vì em thấy chiếc bút chì này thơm và màu đẹp thôi mà, chị Kiri.''

Kiri nghiêng đầu ghi chép vào cuốn sổ tay nhỏ: ``Ghi nhận: hành vi của con người bị chi phối bởi cảm xúc thẩm mỹ nhiều hơn tư duy tối ưu hóa.''

Từ mặt đồng hồ trên tay tôi, luồng sáng xanh của Game khẽ nhấp nháy, âm thanh máy móc vang lên có phần bất lực: ``\eng{Kuon-user, tôi thấy Kiri đã truy vấn hệ thống dữ liệu của ta tổng cộng bốn mươi bảy câu hỏi liên quan đến thuật toán lượng tử trong vòng ba mươi phút vừa qua rồi. Bộ vi xử lý của tôi đang có dấu hiệu quá nhiệt nhẹ vì phải giải thích các khái niệm triết học cho một cô hầu gái rồi đấy.}''

Tôi bật cười, gõ nhẹ vào mặt đồng hồ: ``\eng{Ông luôn tự hào mình thông thái nhất thế gian cơ mà. Giờ gặp đối thủ xứng tầm rồi đấy, ráng mà trả lời cho em nó đi chứ.}''

\ct

Không chỉ gắn kết với các thành viên ruột thịt nhà Hikawa, sự xuất hiện của ba mẹ con còn mang đến những làn gió mới cho những người bạn đang cư ngụ tại các tầng trên của ký túc xá.

Tại tầng sáu, Kson bước ra khỏi phòng với bộ đồ thể thao rộng thùng thình và một lon nước ngọt trên tay. Nàng Tổng trưởng khoanh tay tựa lưng vào lan can, thích thú ngắm nhìn Toa đang một tay cắp hai chiếc lốp xe máy đi phăm phăm qua hành lang: ``Ồ, Kuon-paisen tuyển đâu ra cô em hầu gái xịn sò thế? Sức mạnh này khéo vặn gãy cổ mấy tên du côn đường phố chỉ trong một nốt nhạc đấy chứ.''

Kson bước tới, vỗ vai Toa cười lớn: ``Này, hôm nào rảnh em dạy cho vài câu tiếng lóng đi săn đêm nhé, bảo đảm ra đường không thằng nào dám nhìn đểu luôn!''

Tôi đang đứng ở chân cầu thang nghe thấy thế vội vàng lao tới xua tay rối rít: ``Xin em đấy. Đừng có dạy hư cún nhà anh, con bé còn ngây thơ lắm.''

Kson cười sảng khoái, tiếng cười vang vọng khắp cả dãy hành lang tầng sáu.

Ở tầng hai, Delutaya vốn là một cô bé khép kín và nhút nhát vì những tổn thương trong quá khứ. Ban đầu, khi thấy ba người lạ mặt có tai và đuôi thú xuất hiện, nàng Tam giá hơi rụt rè nép sau cánh cửa phòng. Thế nhưng, sự chân thành và ngây thơ thuần khiết của ba mẹ con đã nhanh chóng phá vỡ bức tường ngăn cách ấy.

Nhìn thấy mái tóc xanh lá cây cùng bộ trang phục mà Delutaya hay diện khi đi diễn trên live, Toa tròn mắt thán phục thốt lên: ``Oa, nhìn chị như một ngôi sao nhạc rock vậy đó!''

Delutaya đỏ bừng mặt vì ngượng, hai tay vân vê lọn tóc màu hồng nhạt. Cùng lúc đó, Miyako nhẹ nhàng bước tới, trên tay là một đĩa bánh nướng xốp mật ong thơm phức vừa mới ra lò: ``Delutaya-sama, em vừa nướng mẻ bánh này, hy vọng hợp khẩu vị của người ạ.''

Sự ấm áp, dịu dàng toát ra từ ánh mắt của người mẹ Miyako đã khiến sự phòng thủ trong lòng cô bé cựu Gen 5 tan chảy hoàn toàn. Delu khẽ gật đầu, nở một nụ cười rụt rè nhưng chan chứa niềm vui: ``Em\dots\ em cảm ơn chị, Miyako-san. Chiều nay rảnh, em sẽ chỉ cho Toa và Kiri cách thắt nơ kiểu mới nhé.''

Thế nhưng, cuộc gặp gỡ kịch tính và dở khóc dở cười nhất trong những ngày đầu ấy lại thuộc về màn đụng độ giữa ba mẹ con tộc cẩu và hai chị em miêu nữ ở tầng tám: Mikeneko và Koneko.

Đó là vào một buổi chiều muộn, khi Mikeneko cùng cô em gái Koneko từ trên phòng bước xuống khu vực phòng khách tầng năm để lấy nước giải khát.

Vừa bước qua ngưỡng cửa, chiếc mũi nhỏ nhắn của nàng Mèo hồng bỗng giật giật hai cái. Nàng miêu nữ mang phong cách Gothic Lolita với mái tóc hồng ngắn và đôi mắt đỏ lập tức khựng lại, đồng tử co rút thành một đường thẳng đứng, toàn bộ lông tơ trên đôi tai mèo dựng đứng lên vì cảnh giác: ``Khoan đã\dots\ Mùi gì thế này? Mùi này\dots\ nồng nặc mùi của loài chó!''

Đúng lúc đó, Toa từ trong gian bếp hớn hở bước ra, trên tay cầm chiếc khăn lau bàn.

Nhìn thấy đôi tai mèo màu hồng trên đầu Mikeneko đang giật giật liên hồi, bản năng săn mồi và sự tò mò của loài chó trong người Toa bỗng chốc trỗi dậy. Cô bé dừng phắt lại, nghiêng đầu sang một bên, hai tai vểnh cao hết cỡ, mắt dán chặt vào đôi tai mèo trước mặt.

Và rồi, hoàn toàn theo phản xạ tự nhiên của một chú chó khi gặp mèo lạ, Toa há miệng, cất giọng thật to và dõng dạc: ``Gâu!''

Tiếng sủa bất ngờ vang lên làm rung chuyển cả gian phòng khách.

Mikeneko giật bắn mình hét toáng lên một tiếng thất thanh, cả người co rúm lại rồi nhảy phắt lên thành tựa lưng của chiếc ghế sofa da, hai chân co lên ngực, hai tay ôm chặt lấy chiếc gối tựa: ``Áaaa! Nó vừa sủa! Nó dám sủa vào mặt mình! Kuon-senpai mau ra đây, đuổi nó ra ngoài ngay, em ghét chó nhất trên đời này!''

Toa đứng ngơ ngác, chớp chớp mắt không hiểu vì sao người ta lại sợ hãi đến thế.

Thấy chị gái hoảng loạn, cô bé mèo nhỏ Koneko đứng bên dưới lại không hề tỏ ra sợ hãi. Ngược lại, cô bé tò mò chầm chậm tiến lại gần Kiri đang đứng bên cạnh tủ sách. Cô bé mèo khẽ hếch mũi ngửi ngửi mùi hương hoa cúc thanh nhẹ trên tạp dề của Kiri, đôi mắt tròn xoe ánh lên vẻ thích thú.

Kiri không hề tỏ ra bối rối. Cô bé từ tốn quỳ một chân xuống sàn nhà để ngang bằng tầm mắt với Koneko, hai tay cẩn thận nâng chiếc khay bạc, trên đó là một chiếc đĩa sứ nhỏ chứa chiếc bánh quy hình chú cá nướng giòn rụm: ``Xin chào Koneko-sama. Đây là bánh quy cá ngừ do mẹ em nướng, xin phép được gửi tặng người làm quà chào hỏi ạ.''

Koneko nhìn chiếc bánh cá ngừ thơm lừng, hai mắt sáng rực lên. Cô bé nhận lấy chiếc bánh, cắn một miếng nhỏ rồi khẽ kêu lên một tiếng ``Meo'' đầy mãn nguyện, chiếc đuôi mèo ngoe nguẩy vòng quanh chân Kiri.

Cùng lúc đó, Miyako từ tốn bước ra từ gian bếp, trên tay bưng một chiếc khay gỗ đựng ấm trà hoa cúc mật ong bốc khói nghi ngút. Cô bước đến bên chiếc ghế sofa nơi Mikeneko vẫn đang ngồi thu lu với vẻ mặt ấm ức, khẽ cúi người thật sâu, giọng nói dịu dàng như làn nước mát mùa thu: ``Thành thật xin lỗi Mikeneko-sama. Toa tính tình còn trẻ con và xốc nổi, đã làm kinh động đến người. Em có pha một ấm trà hoa cúc mật ong giúp an thần và làm dịu tinh thần, kính mời người dùng thử để hạ hỏa ạ.''

Mùi trà hoa cúc thơm ngát hòa cùng mật ong ngọt ngào lan tỏa trong không khí, kết hợp với phong thái nhã nhặn, cung kính tuyệt đối của Miyako đã khiến cơn giận dỗi của Mikeneko dần xẹp xuống.

Nàng miêu nữ liếc nhìn đĩa bánh cá của em gái Koneko đang ăn ngon lành, lại nhìn tách trà thơm phức trước mặt, hai má khẽ ửng hồng. Cô nàng chầm chậm tụt xuống ghế, đón lấy tách trà, lí nhí trong miệng: ``Hừm\dots\ Coi như chị biết điều đấy. Lần sau bảo con bé tai dài kia cấm có được sủa bất thình lình như thế nữa, đau tim chết đi được\dots''

Toa ở phía sau vội vàng gãi đầu, cúi gập người xin lỗi: ``Em xin lỗi, Mikeneko-sama! Tại em thấy tai chị giật giật dễ thương quá nên lỡ mồm thôi mà!''

Câu khen ngợi ngây ngô của Toa khiến nàng Mèo hồng khựng lại một nhịp, quay mặt đi chỗ khác giả vờ nhấp trà để giấu đi nụ cười ngượng ngùng. Trận chiến truyền kiếp giữa mèo và chó tưởng chừng sẽ nổ ra nảy lửa, rốt cuộc lại được hóa giải một cách êm đẹp và ngọt ngào như thế.

\ct

Đêm hôm đó, đồng hồ điểm mười một giờ khuya.

Cả ký túc xá đã chìm vào giấc ngủ yên tĩnh. Tôi vẫn ngồi bên bàn làm việc ở tầng năm, ánh đèn bàn màu vàng ấm soi rọi lên những tập tài liệu quy hoạch sự kiện của công ty đang xem dở.

*CỐC CỐC CỐC* Một tiếng gõ cửa khẽ khàng vang lên ba nhịp.

``Vào đi,'' tôi cất tiếng.

Cánh cửa gỗ mở ra thật êm. Miyako bước vào, trên tay bưng một tách trà thảo mộc còn nghi ngút khói. Cô đã cởi bỏ chiếc bờm cài đầu và chiếc tạp dề trắng, chỉ còn mặc chiếc đầm đen tay dài trang nhã.

Miyako bước nhẹ nhàng tới mép bàn, cẩn thận đặt tách trà xuống chiếc lót ly bằng gỗ trước mặt tôi: ``Cậu chủ, đã khuya rồi, xin mời Cậu chủ dùng một tách trà ấm để thư giãn tinh thần ạ.''

Tôi ngước nhìn cô, mỉm cười đón lấy tách trà: ``Cảm ơn em, Miyako. Những ngày vừa rồi vất vả cho ba mẹ con quá.''

Miyako lắc đầu nhẹ nhàng, đôi mắt vàng kim long lanh dưới ánh đèn bàn: ``Không hề vất vả chút nào đâu ạ. Được làm việc, được phục vụ cho Cậu chủ và đại gia đình là niềm hạnh phúc lớn lao nhất mà ba mẹ con chúng em có được sau những ngày tháng tăm tối đã qua.''

Tôi nhấp một ngụm trà ấm, cảm nhận vị ngọt thanh đọng lại nơi đầu lưỡi. Nhìn dáng vẻ khép nép lễ phép của cô, tôi khẽ thở dài một tiếng rồi nói bằng giọng chân thành: ``Miyako này, anh đã nói mấy lần rồi đấy. Ở ngôi nhà này, mọi người sống với nhau như ruột thịt, không có chuyện chủ tớ phong kiến như trong mấy bộ phim truyền hình cổ trang đâu. Em không cần phải lúc nào cũng xưng hô là `Cậu chủ' với anh, cũng không cần phải gọi Kaiho là `Đại tiểu thư' hay gọi mấy đứa em gái anh là `Tiểu thư' này nọ. Cứ gọi tên bình thường như `anh Kuon', `chị Kaiho', `Suzuno' là được rồi mà.''

Miyako ngẩn người ra một thoáng. Đôi tai cún trắng muốt trên đầu cô khẽ run rẩy nhè nhẹ. Cô đưa bàn tay lên áp chặt trước ngực áo, nơi chiếc nơ vàng đồng vẫn được cài trang trọng, ánh mắt nhìn tôi đong đầy sự kiên định cùng lòng tôn kính sâu sắc: ``Dạ\dots\ em hiểu tấm lòng bao dung của Cậu chủ. Nhưng xin Cậu chủ hãy cho phép em được gọi như thế ạ. Với ba mẹ con chúng em, việc được tôn xưng Cậu chủ và các Đại tiểu thư không phải là sự gò bó, mà là biểu tượng cho danh dự, cho sự trung thành tuyệt đối và là sợi dây kết nối thiêng liêng nhất mà chúng em gìn giữ đối với người chủ nhân đã ban cho chúng em sinh mệnh thứ hai này. Nếu bắt em phải xưng hô suồng sã ngang hàng, trong lòng em sẽ cảm thấy bứt rứt và áy náy khôn nguôi.''

Nhìn ánh mắt rực rỡ sự kiên định và tấm lòng chân thành không chút toan tính của người phụ nữ mang linh hồn của người bạn bốn chân năm nào, tôi khẽ mỉm cười, đưa tay xoa nhẹ gáy rồi buông một tiếng thở phào chịu thua:  ``Thôi thì\dots\ gọi theo cách nào mà em cảm thấy thoải mái và tự nhiên nhất là được.''

Miyako mỉm cười rạng rỡ, khóe mắt ánh lên những tia sáng ấm áp, khẽ cúi người thật sâu: ``Em xin cảm tạ Cậu chủ. Chúc Cậu chủ có một giấc ngủ thật ngon ạ.''

Cô xoay người, nhẹ nhàng rón rén bước từng bước êm ái ra phía cửa rồi khép lại cẩn thận.

Tôi nhìn theo bóng lưng của cô, lòng bỗng dâng lên một cảm giác ấm áp lạ thường.

Từ mặt chiếc đồng hồ thông minh trên cổ tay, Game khẽ nhấp nháy luồng sáng xanh lục, âm thanh máy móc vang lên phẳng lặng nhưng mang theo chút hài lòng hiếm thấy: ``Kuon-user, theo dữ liệu quét nhịp tim và sóng não từ xa, chỉ số hạnh phúc và mức độ gắn kết tâm lý của ba thực thể sinh học đó đối với gia đình này đạt mức chín mươi chín phẩy tám phần trăm. Đồng thời, sự hiện diện của họ đã giúp giảm tải ba mươi tám phần trăm mức độ căng thẳng chung của cả ký túc xá. Tôi phải thừa nhận\dots\ quyết định giữ họ lại của cậu không hề tồi.''

Tôi khẽ bật cười, tựa lưng vào thành ghế, nhấp thêm một ngụm trà ấm rồi nhìn ra màn đêm lung linh ánh đèn của thành phố Kawachi qua khung cửa sổ.

Những ngày mới đã thực sự bắt đầu, và tôi biết, ngôi nhà này từ nay sẽ không bao giờ còn phải chịu đựng sự cô đơn hay những khoảng trống lạnh lẽo thêm một lần nào nữa.

\end{document}